% ====================================================================
\section{Setup and procedure}\label{sec:setup}
% ====================================================================

\begin{figure}[t]
  \centering
  \resizebox{\textwidth}{!}{\begin{tikzpicture}[
  font=\footnotesize,
  box/.style={draw, rounded corners=1.5pt, minimum height=7.5mm, minimum width=12mm, align=center, fill=white, inner sep=2pt},
  mw/.style={-{Latex[length=2mm]}, line width=1.1pt, draw=black!80},
  mwb/.style={{Latex[length=2mm]}-{Latex[length=2mm]}, line width=1.1pt, draw=black!80},
  lf/.style={-{Latex[length=1.6mm]}, densely dashed, draw=black!65},
  lbl/.style={font=\scriptsize, text=black!75}]
  % ---------------- Microwave bridge
  \node[box] (dro)  at (-8.0,0) {DRO\\\SI{9.36}{GHz}};
  \node[box] (iso)  at (-6.2,0) {Isolator};
  \node[box] (pick) at (-4.3,0) {Directional\\coupler};
  \node[box] (att1) at (-2.3,0) {X375A\\0--20\,dB};
  \node[box, minimum width=12mm, minimum height=12mm] (tee) at (0,0) {Magic-T};
  \node[box] (att2) at (2.3,0) {X382A\\0--50\,dB};
  \node[box] (plunger) at (4.3,0) {Plunger\\$z$};
  \node[box] (counter) at (-4.3,1.55) {Counter\\HP\,5342A};
  \node[box] (diode) at (0,1.75) {Diode\\HP\,8471E};
  \node[box, minimum width=13mm] (cav) at (0,-1.85) {Cavity\\TE$_{102}$};
  \fill[black!12, draw=black!70] (-1.35,-2.55) rectangle (-0.95,-1.15);
  \fill[black!12, draw=black!70] ( 0.95,-2.55) rectangle ( 1.35,-1.15);
  \node[lbl, align=left, anchor=west] at (1.5,-1.85) {Magnet with\\modulation coils};
  \draw[mw] (dro) -- (iso);
  \draw[mw] (iso) -- (pick);
  \draw[mw] (pick) -- (att1);
  \draw[mw] (pick) -- (counter);
  \draw[mw] (att1) -- node[above, lbl] {$\Sigma$} (tee);
  \draw[mwb] (tee) -- (att2);
  \draw[mwb] (att2) -- (plunger);
  \draw[mwb] (tee) -- (cav);
  \draw[mw] (tee) -- node[right, lbl] {$\Delta$} (diode);
  % ---------------- Electronics
  \node[box, minimum width=27mm, minimum height=10mm] (lockin) at (6.3,1.75) {Lock-in SR830};
  \node[box] (pc) at (6.3,3.25) {PC, LabVIEW};
  \node[box] (delta) at (3.4,-3.9) {Power supply\\SM400};
  \node[box] (modamp) at (6.3,-3.9) {Mod.\\amplifier};
  \node[box] (gauss) at (-3.4,-3.9) {Gaussmeter\\F.\,W.\,Bell 5080};
  \draw[lf] (diode) -- node[above, lbl, pos=0.5] {In A} (lockin.west |- diode);
  \draw[lf] (pc) -- node[right, lbl] {GPIB} (lockin);
  \draw[lf] (5.3,1.25) -- node[pos=0.12, left, lbl] {Aux Out} (5.3,-3.9) -- (delta.east);
  \draw[lf] (7.3,1.25) -- node[pos=0.12, right, lbl] {Sine Out, \SI{133}{Hz}} (7.3,-3.9) -- (modamp.east);
  \draw[lf] (delta.west) -| node[pos=0.75, right, lbl] {$I$} (1.15,-2.55);
  \draw[lf] (modamp.south) |- (0,-4.75) -| (-1.15,-2.55);
  \draw[lf] (gauss.north) |- node[pos=0.75, above, lbl] {Hall probe} (-1.35,-1.85);
  \draw[lf] (gauss.west) -| (-5.25,2.6) -| node[pos=0.25, above, lbl] {Aux In} (5.95,2.25);
\end{tikzpicture}
}
  \caption{X-band reflection spectrometer, drawn with the arms of the magic T as on the real T-piece. Microwave path
  (thick): the DRO source feeds the main attenuator X375A through an isolator and a directional coupler (frequency
  counter). Placed \emph{before} the magic T, it attenuates the cavity and reference waves together. The reference arm
  (top, opposite the cavity) with the reference-arm attenuator X382A and the plunger sets the magnitude and phase of the
  reference. The diode sits at the $\Delta$ port (right). Signal and control lines (dashed): the lock-in supplies the
  control voltage $U$ for the magnet power supply (Aux Out, a DAC output) and the modulation (Sine Out, \qty{133}{Hz})
  and reads the Hall probe (Aux In, an ADC input) and the diode (In A).}
  \label{fig:setup}
\end{figure}

\paragraph{Microwaves.}
\Cref{fig:setup} shows the setup, and \cref{tab:settings} lists all instrument settings. The source is a tunable
dielectric resonator oscillator (DRO, $\geq\qty{13}{dBm}$, i.e.\ at least \qty{20}{mW}). A directional coupler sends
part of its power to the frequency counter. The empty cavity (TE$_{102}$) resonates at \qty{9.3875}{GHz} with a
width of \qty{10.1}{MHz} (network analyzer), so $\QL\approx930$. After every sample change I set X375A to strong
attenuation (\qty{20}{dB}) and retuned the source to the \enquote{dip}, the minimum of the diode voltage where the
cavity absorbs the power. Then I set X375A to \qty{2}{dB}, X382A to \qtyrange{0}{1}{dB} and the plunger to
\qty{20.00}{mm}. With a sample, the resonance was at \qty{9.3594}{GHz} (DPPH) or \qty{9.3527}{GHz} (\ce{MnCl2}). The
detector diode (HP\,8471E) sits at the difference port ($\Delta$) of the magic T.

\paragraph{Fields.}
The electromagnet runs on a current-regulated power supply whose set point is the control voltage
$U=\qtyrange{0}{8}{V}$ from the lock-in. At the DPPH resonance the current is \qty{0.43}{A}. The lock-in's Sine Out
(\qty{133}{Hz}) drives the modulation coils via an amplifier. A gaussmeter (F.\,W.\,Bell\,5080) measures the field
with a Hall probe. I recorded its analog output (\qty{1}{V} corresponds to \qty{1}{T}) and read its display directly. A LabVIEW
program sweeps $U$ in small steps and stores the lock-in outputs $X$ and $Y$ at each step. I set the lock-in phase so
that practically the whole signal lies in $X$ ($Y/X=\num{-0.05}$), and I analyze this component as $S=-X$. I set the
phase with Auto Phase at
$U=\qty{5.9}{V}$, on the high-field flank of the \ce{MnCl2} line (\cref{tab:settings}). There the derivative of the
absorption is negative, but Auto Phase makes the signal positive, so the lock-in phase is \ang{180} off the usual
convention and the raw $X$ has the opposite sign. All fits, figures and signed amplitudes therefore refer to $S=-X$. With it, the derivative of an
absorption line is positive on the rising (low-field) side and negative on the falling (high-field) side, as in
\cref{fig:theory}c.

\paragraph{Measurement program.}
\Cref{tab:scans} lists the analyzed measurements. Before each DPPH scan I set $U$ from \qty{0}{V} to the start value
and held it for \qtyrange{30}{60}{s}. A typical sweep ran at \qty{0.48}{mV/s} of control voltage. To make this independent of the equipment, I convert it
with the slope $\dd B/\dd U=\qty{60.76}{mT/V}$ at the DPPH line (\cref{eq:fieldscale}) and the coil current of
\qty{0.43}{A} at $U=\qty{4.95}{V}$, i.e., \qty{87}{mA/V} if the current is proportional to the set point. The field
then rose at \qty{29}{\micro T/s} (\qty{0.29}{G/s}) and the coil current at about
\qty{42}{\micro A/s}. The DPPH line ($\Bpp=\qty{0.24}{mT}$, \qty{4}{mV} of control voltage) was crossed in about
\qty{8}{s}. This is $27\tau$ for the lock-in time constant $\tau=\qty{300}{ms}$ (the averaging time of its low-pass
filter) and satisfies the rule of thumb from \cref{sec:modulation}. \Cref{tab:scans} gives the field rate of every
scan. These rates use the static slope; during a sweep the field lags behind the set value (\cref{sec:field}), so the
actual rate may be up to \qty{22}{\percent} lower. The DPPH sample is a large
volume of powder. I left out the cut-off scans 07 and 17 (repeated as 07b and 17b) and the first point of every file,
where the signal had not settled. The DPPH search scan 06 overloaded the lock-in, so I only use its zero
crossing, noted in my lab notebook (\cref{sec:mncl2}).

\begin{table}[t]
  \centering\footnotesize
  \caption{Analyzed measurements (23 Oct 2026). $v$: sweep rate of the control voltage; $\dot B$: the corresponding
  field rate $v\,\dd B/\dd U$ with the static slope at the line (\qty{60.76}{mT/V} for DPPH, \qty{61.1}{mT/V} for
  \ce{MnCl2} and \ce{Cu^2+}; for 01--03 the mean slope over the full range). Mod.: Sine Out modulation
  voltage, in parentheses the effective value after amplifier limiting (\cref{sec:modamp}). Attenuators
  main attenuator X375A (cavity and reference) and reference-arm attenuator X382A (reference only) in dB. Sens.: lock-in sensitivity (full scale) in \si{\micro V}.}
  \label{tab:scans}
  \setlength{\tabcolsep}{2.6pt}
  \begin{tabularx}{\textwidth}{@{}l l l l S[table-format=-2.2] S[table-format=-4.0] l c c c >{\raggedright\arraybackslash}X@{}}
    \toprule
    No. & Time & Sample & $U$ (V) & {$v$ (\si{mV/s})} & {$\dot B$ (\si{\micro T/s})} & Mod.\ (V) & X375A & X382A & Sens. & Purpose \\
    \midrule
    01--03 & 10:00 & empty & \numrange{0.10}{8.05} & 25 & 1500 & -- & -- & -- & -- & Hall probe up/down/up \\
    22 & 11:41 & \ce{MnCl2} & \numrange{2.00}{7.98} & 2.86 & 175 & 4 (1.8) & 2 & 1 & 2 & \ce{MnCl2} \\
    07b, 08 & 13:33 & DPPH & \numrange{4.84}{5.13} & 0.95 & 58 & 3; 2 (1.8) & 2 & 1 & 200 & modulation \\
    09--12 & 13:56 & DPPH & \numrange{4.90}{5.01} & 0.48 & 29 & 1--0.1 & 2 & 1 & 200--10 & modulation \\
    13--17b & 14:27 & DPPH & \numrange{4.91}{4.99} & 0.48 & 29 & 1 & 2 & 0--20 & 100--20 & reference arm \\
    18 & 15:14 & DPPH & \numrange{4.99}{4.91} & -0.48 & -29 & 1 & 2 & 6 & 50 & down-sweep \\
    24--27 & 15:24 & DPPH & \numrange{4.91}{4.99} & 0.48 & 29 & 1 & 6--20 & 0 & 100--5 & power series \\
    20 & 15:55 & DPPH & \numrange{4.91}{4.99} & 0.48 & 29 & 1 & 2 & 0 & 200 & dispersion, plunger \qty{25.60}{mm} \\
    19 & 16:05 & -- & \numrange{4.90}{5.00} & & & -- & -- & -- & -- & gaussmeter display at sample \\
    29 & 16:58 & \ce{Cu^2+} & \numrange{4.40}{5.19} & 4.76 & 292 & 4 (1.8) & 2 & 0 & 2 & search \\
    \bottomrule
  \end{tabularx}
\end{table}

\begin{table}[t]
  \centering\footnotesize
  \caption{Instrument settings. Values that vary between scans are in \cref{tab:scans}.}
  \label{tab:settings}
  \begin{tabularx}{\textwidth}{@{}>{\raggedright\arraybackslash}p{0.25\textwidth} >{\raggedright\arraybackslash}X@{}}
    \toprule
    Instrument & Setting \\
    \midrule
    Microwave source & Tunable dielectric resonator oscillator (DRO), $\geq\qty{13}{dBm}$; tuned to the cavity dip after
      every sample change, with X375A at \qty{20}{dB} while tuning \\
    Frequency counter & Via a directional coupler: \qty{9.3594}{GHz} (DPPH), \qty{9.3527}{GHz} (\ce{MnCl2}),
      \qty{9.199}{GHz} (\ce{Cu^2+}, detuned) \\
    Cavity & Rectangular, TE$_{102}$; empty: resonance \qty{9.3875}{GHz}, width \qty{10.1}{MHz} ($\QL\approx930$) \\
    Attenuator X375A (main) & \qty{2}{dB}; \qtyrange{6}{20}{dB} in the power series 24--27 \\
    Attenuator X382A (reference arm) & \qty{1}{dB} (scans 07b--12, 22) or \qty{0}{dB} (13, 20, 24--29); \qtyrange{0}{20}{dB}
      in 13--17b \\
    Plunger (reference arm) & \qty{20.00}{mm} (absorption); \qty{25.60}{mm} (dispersion, scan 20) \\
    Detector & Diode HP\,8471E at the $\Delta$ port of the magic T, connected to lock-in input A and, via a T-piece, to an
      oscilloscope \\
    Lock-in amplifier & Stanford Research SR830, internal reference \qty{133}{Hz}; time constant \qty{300}{ms}
      (\qty{1}{s} in scan 22), filter slope \qty{12}{dB/oct}, reserve Low Noise; phase set once with Auto Phase at
      $U=\qty{5.9}{V}$ before scan 22 and left unchanged afterwards; sensitivity in \cref{tab:scans} \\
    Field modulation & SR830 Sine Out, \qtyrange{0.1}{4}{V} (rms), drives the modulation coils via an amplifier, which
      limits above about \qty{1.8}{V} \\
    Magnet power supply & Delta SM400, current-controlled; set point $U=\qtyrange{0}{8}{V}$ from SR830 Aux Out (DAC);
      \qty{0.43}{A} at the DPPH line \\
    Gaussmeter & F.\,W.\,Bell 5080 with Hall probe (ranges \qtylist{30;300;3000}{mT}, bandwidth dc to
      \qty{20}{kHz}~\cite{fwbell}), range setting \enquote{0.000\,T}; analog output (\qty{1}{V}
      corresponds to \qty{1}{T}) recorded via SR830 Aux In (ADC), display read directly \\
    Field sweep (LabVIEW program EPR\_1) & DPPH: steps of \qty{1}{mV} (\qty{2}{mV} in 07b, 08), \qty{2.1}{s} per point;
      \ce{MnCl2} (22): steps of \qty{20}{mV}, \qty{7}{s} per point; before each scan $U$ set from \qty{0}{V} to the start
      value and held for \qtyrange{30}{60}{s}; back to \qty{0}{V} after each scan \\
    Hall-probe scans 01--03 & LabVIEW program Main\_Field\_Calibration\_2: from \qty{0.1}{V} in 158 steps of \qty{50}{mV},
      \qty{1}{s} wait per step, probe reading averaged over $20\times\qty{50}{ms}$ \\
    \bottomrule
  \end{tabularx}
\end{table}
